\section{实验计划与待完成实验}
\label{sec:experiments-tbd}

本节区分已完成实验与待完成实验。当前完成的 E11 局部几何实验见第~\ref{sec:e11} 节。下面列出的实验尚未完成，本文暂时保留空位。

\subsection{实验一：Activation perturbation diagnostics}

\paragraph{状态.} \TBD。

\paragraph{目标.} 直接测量 SpecGrad / polar update 是否在 activation perturbation 上表现出 operator-norm 意义下的稳定性。

对每层记录
\[
\Delta Z_{\ell,\mathrm{direct}}
=(W_\ell^{t+1}-W_\ell^t)A_{\ell-1}^t.
\]
测量：
\[
\frac{\|\Delta Z_{\ell,\mathrm{direct}}\|_F}{\|Z_\ell^t\|_F},
\qquad
\frac{\|\Delta Z_{\ell,\mathrm{direct}}\|_{\op}}{\|Z_\ell^t\|_{\op}},
\]
以及 per-sample 最大扰动：
\[
\max_b
\frac{\|(W_\ell^{t+1}-W_\ell^t)h_{\ell-1}(x_b)\|_2}
{\|h_{\ell-1}(x_b)\|_2}.
\]

\paragraph{预期.} SpecGrad 应更自然地控制 operator/per-sample 指标，而不一定降低 batch-Frobenius movement。

\paragraph{结果.} \TBD。

\subsection{实验二：Activation stable rank 与层级条件}

\paragraph{状态.} \TBD。

计算
\[
\srank(A_{\ell-1})=\frac{\|A_{\ell-1}\|_F^2}{\|A_{\ell-1}\|_{\op}^2},
\qquad
\nrank(G_\ell)=\frac{\|G_\ell\|_*^2}{\|G_\ell\|_F^2}.
\]
检验条件
\[
\nrank(G_\ell)>\srank(A_{\ell-1})
\]
是否预测该层 polar/SpecGrad 的一阶优势。

\paragraph{结果.} \TBD。

\subsection{实验三：Downstream-aware activation perturbation}

\paragraph{状态.} \TBD。

目标是测量
\[
\|B_\ell D_\ell A_{\ell-1}\|_F
\]
或使用 JVP 估计
\[
\|J_\theta D_\ell\|_F.
\]
比较
\[
D_{\mathrm{spec}}=\operatorname{polar}(G_\ell)
\]
与
\[
D_{\mathrm{gd}}=G_\ell/\|G_\ell\|_F.
\]
核心指标为
\[
\frac{\langle G_\ell,D_{\mathrm{spec}}\rangle^2/\|J_\theta D_{\mathrm{spec}}\|_F^2}
{\langle G_\ell,D_{\mathrm{gd}}\rangle^2/\|J_\theta D_{\mathrm{gd}}\|_F^2}.
\]

\paragraph{结果.} \TBD。

\subsection{实验四：长尾小 batch 激活实验}

\paragraph{状态.} \TBD。

构造长尾分类或 synthetic long-tail mode 数据。检验 rare direction 可见性条件
\[
\frac{Bp_c}{1-\beta}\gtrsim 1.
\]
记录 tail-class loss、tail activation drift、tail activation coverage、minibatch spectral-tail consistency 与 effective nuclear rank。

\paragraph{结果.} \TBD。

\subsection{实验五：现代架构激活扰动 benchmark}

\paragraph{状态.} \TBD。

在 Transformer、现代 MLP 或 CNN 架构中测量
\[
\|\Delta A_\ell\|_F,
\qquad
\|\Delta A_\ell\|_{\op},
\qquad
\srank(A_\ell),
\qquad
\|J_\theta D_\ell\|_F.
\]
本实验不以 attention 为核心，而是直接观察 activation perturbation。

\paragraph{结果.} \TBD。
